\section{Scanner Hand-Written dalam C}
\label{sec:ch07-tangan}

\subsection{Struktur Data}

Lexer tangan menggunakan struktur \code{Lexer} yang menyimpan sumber,
posisi, baris, kolom, dan status kesalahan. Struktur ini didefinisikan
pada \file{lexer.h}, seperti ditunjukkan pada
\lstref{lst:ch07-lexer-h}.

\includecode[
  language=C,
  firstline=1,
  lastline=25,
  firstnumber=1,
  caption={Antarmuka \file{lexer.h}},
  label={lst:ch07-lexer-h}
]{../../code/subsetc/include/lexer.h}

Pemisahan header dan implementasi memudahkan pengujian dan
memungkinkan driver lain memakai lexer tanpa mengetahui detail
internal.

\subsection{Operasi Dasar}

Tiga operasi dasar menjadi fondasi scanner:
\begin{itemize}[leftmargin=*]
  \item \code{peek()} mengembalikan karakter saat ini tanpa memajukan
        posisi;
  \item \code{peek\_next()} mengembalikan karakter berikutnya;
  \item \code{advance()} memajukan posisi dan memperbarui baris serta
        kolom.
\end{itemize}

Ketiga operasi ini menangani batas akhir masukan dengan
mengembalikan \code{EOF}. Semua operasi lain dibangun di atas
ketiganya, sehingga logika pelacakan lokasi hanya ditulis sekali.

\subsection{Melewati Spasi dan Komentar}

Fungsi \code{skip\_space} membuang spasi, tab, baris baru, komentar
baris, dan komentar blok. Untuk komentar blok yang tidak ditutup,
fungsi melaporkan kesalahan dan menandai lexer berada dalam keadaan
error. Potongan implementasinya ditunjukkan pada
\lstref{lst:ch07-skip-space}.

\includecode[
  language=C,
  firstline=51,
  lastline=95,
  firstnumber=51,
  caption={Fungsi \code{skip\_space} pada \file{lexer.c}},
  label={lst:ch07-skip-space}
]{../../code/subsetc/src/lexer.c}

Perhatikan bahwa komentar tidak menghasilkan token. Lexer hanya
memperbarui posisi dan melanjutkan ke karakter berikutnya.

\subsection{Fungsi Utama \code{lexer\_next}}

Fungsi \code{lexer\_next} adalah inti scanner. Alurnya:
\begin{enumerate}[label=\arabic*.,leftmargin=*]
  \item lewati spasi dan komentar;
  \item catat baris, kolom, dan posisi awal;
  \item jika akhir masukan, kembalikan \code{TOK\_EOF};
  \item jika huruf atau garis bawah, pindai identifier lalu periksa
        tabel kata kunci;
  \item jika digit, pindai bilangan bulat atau pecahan;
  \item jika kutip, pindai char atau string;
  \item jika operator, periksa operator dua karakter lalu satu
        karakter;
  \item jika tidak dikenal, hasilkan \code{TOK\_ERROR}.
\end{enumerate}

Bagian awal fungsi, termasuk penanganan identifier, ditunjukkan pada
\lstref{lst:ch07-lexer-next-id}.

\includecode[
  language=C,
  firstline=149,
  lastline=177,
  firstnumber=149,
  caption={Awal \code{lexer\_next}: EOF dan identifier},
  label={lst:ch07-lexer-next-id}
]{../../code/subsetc/src/lexer.c}

Penanganan bilangan memisahkan integer dan float berdasarkan ada
tidaknya titik yang diikuti digit. Jika titik tidak diikuti digit,
titik tidak dikonsumsi sebagai bagian bilangan, sehingga dapat
dilaporkan sebagai kesalahan pada langkah berikutnya.

\subsection{Operator dan Tanda Baca}

Operator dua karakter diperiksa lebih dahulu. Potongan kode untuk
operator dan penanganan kesalahan ditunjukkan pada
\lstref{lst:ch07-operator}.

\includecode[
  language=C,
  firstline=274,
  lastline=341,
  firstnumber=274,
  caption={Penanganan operator dan karakter ilegal},
  label={lst:ch07-operator}
]{../../code/subsetc/src/lexer.c}

Setiap cabang mengembalikan token dengan panjang yang tepat. Karena
\code{advance()} sudah dipanggil sebelum \code{switch}, panjang token
mencakup karakter operator. Untuk operator dua karakter,
\code{advance()} dipanggil sekali lagi.

\subsection{Kata Kunci}

Fungsi \code{keyword\_type} membandingkan lexeme dengan tabel kata
kunci. Jika tidak ada yang cocok, lexeme dikembalikan sebagai
\code{TOK\_IDENT}. Pendekatan ini memisahkan daftar kata kunci dari
logika scanning, sehingga menambah kata kunci cukup mengubah tabel.
